% !PRINTED 519
\weberpage{519}
%<<CONT
wird, falls wir den Exponenten gleich Null setzen, wenn einer der Primfactoren in dem betreffenden Functionale nicht aufgeht. Ein gemeinsamer Theiler der Zahlen $\alpha, \beta, \ldots$ kann keine anderen Primfactoren als $\pi_1, \pi_2, \ldots$ enthalten, und jeder gemeinsame Theiler von $\alpha, \beta, \ldots$ hat die Form
\weq{8}{\varepsilon \delta = \pi_1^a \pi_2^b \ldots,}
worin $a$ nicht grösser als die kleinste der Zahlen $a_1, b_1, \ldots$ sein darf, $b$ nicht grösser als die kleinste der Zahlen $a_2, b_2, \ldots$ u. s. f.

Ist $a$ die kleinste unter den Zahlen $a_1, b_1, \ldots$, $b$ die kleinste unter den Zahlen $a_2, b_2, \ldots$, so ist die in (8) dargestellte Zahl $\delta$ der grösste gemeinschaftliche Theiler der Zahlen $\alpha, \beta, \ldots$ In Worten ausgedrückt:

Man erhält den grössten gemeinschaftlichen Theiler mehrerer ganzer Functionale $\alpha, \beta, \ldots$, wenn man ein Product aus Primfactoren bildet, in das man jeden Primfactor so oft aufnimmt, als er in jeder der Zahlen $\alpha, \beta, \ldots$ aufgeht.

Zwei Zahlen sind relativ prim, wenn sie keinen gemeinschaftlichen Primfactor enthalten.

Dem entsprechend definiren wir als das \gesperrt{kleinste gemeinschaftliche Multiplum} $\mu$ der Functionale $\alpha, \beta, \ldots$ ein Product aus Primfactoren, in das wir einen Factor $\pi$ so oft aufnehmen, dass er in keinem der Functionale $\alpha, \beta, \ldots$ öfter als in $\mu$ aufgeht. Dieses Functional $\mu$ hat dann, ebenso wie jedes mit $\mu$ associirte Functional, die doppelte, und, wie wir hinzufügen können, charakteristische Eigenschaft, dass es durch jedes der Functionale $\alpha, \beta, \ldots$ theilbar ist, und dass jedes andere Functional, das durch $\alpha, \beta, \ldots$ theilbar ist, auch durch $\mu$ theilbar ist.

Das kleinste gemeinschaftliche Multiplum zweier relativer Primfunctionale ist ihr Product.

Man kann diese Sätze anwenden, um gebrochene Functionale \gesperrt{in der einfachsten Gestalt oder als reducirte Brüche dar}zustellen, indem man Zähler und Nenner in ihre Primfactoren zerlegt und den grössten gemeinschaftlichen Theiler weghebt. Zähler und Nenner eines reducirten Bruches sind durch den Bruch selbst völlig bestimmt, abgesehen von einem gemeinschaftlichen Einheitsfactor, der unbestimmt bleibt.

Ebenso kann man eine beliebige Zahl gegebener Brüche auf gemeinsamen Nenner, den Hauptnenner, bringen, indem man ein gemeinsames Multiplum aller gegebenen Nenner als gemeinschaftlichen Nenner wählt.
